% Part: first-order-logic
% Chapter: sequent-calculus
% Section: rules-and-proofs

\documentclass[../../../include/open-logic-section]{subfiles}

\begin{document}

\iftag{FOL}
      {\olfileid[id]{fol}{seq}{rul}}
      {\olfileid[id]{pl}{seq}{rul}}

\olsection{Aturan dan \usetoken{P}{derivation}}

Untuk uraian berikut, misalkan $\Gamma, \Delta, \Pi, \Lambda$ menyatakan
barisan berhingga !!{sentence}s.

\begin{defn}[Sekuen]
Sebuah \emph{sekuen} adalah ekspresi berbentuk
\[
\Gamma \Sequent \Delta
\]
dengan $\Gamma$ dan $\Delta$ merupakan barisan berhingga (yang mungkin
kosong) dari !!{sentence}s dalam bahasa $\Lang L$. $\Gamma$ disebut
\emph{anteseden}, sedangkan $\Delta$ disebut \emph{sukseden}.
\end{defn}

\begin{explain}
Gagasan intuitif di balik sebuah sekuen ialah: jika semua !!{sentence}s
dalam anteseden berlaku, maka setidaknya satu !!{sentence} dalam sukseden
berlaku. Artinya, jika $\Gamma = \tuple{!A_1, \dots, !A_m}$ dan
$\Delta = \tuple{!B_1, \dots, !B_n}$, maka $\Gamma \Sequent \Delta$
berlaku jika dan hanya jika
\[
(!A_1 \land \cdots \land !A_m) \lif (!B_1 \lor \cdots \lor
!B_n)
\]
berlaku. Ada dua kasus khusus: ketika $\Gamma$ kosong dan ketika
$\Delta$~kosong. Ketika $\Gamma$ kosong, yaitu $m = 0$, $\quad
\Sequent \Delta$ berlaku jika dan hanya jika $!B_1 \lor \dots \lor !B_n$
berlaku. Ketika $\Delta$ kosong, yaitu $n = 0$, $\Gamma \Sequent \quad$
berlaku jika dan hanya jika $\lnot(!A_1 \land \dots \land !A_m)$ berlaku.
Kita mengatakan bahwa sebuah sekuen valid jika dan hanya jika
!!{sentence} yang bersesuaian dengannya valid.
\end{explain}

Jika $\Gamma$ adalah barisan !!{sentence}s, kita menulis $\Gamma, !A$
untuk hasil menambahkan $!A$ pada ujung kanan~$\Gamma$ (dan $!A, \Gamma$
untuk hasil menambahkan $!A$ pada ujung kiri~$\Gamma$). Jika $\Delta$
juga merupakan barisan !!{sentence}s, maka $\Gamma, \Delta$ adalah
konkatenasi kedua barisan itu.

\begin{defn}[Sekuen Awal]
Sebuah \emph{sekuen awal} adalah sekuen
\iftag{prvFalse,prvTrue}{dengan salah satu bentuk berikut:
  \begin{enumerate}
    \item $!A \Sequent !A$
    \tagitem{prvTrue}{$\quad \Sequent \ltrue$}{}
    \tagitem{prvFalse}{$\lfalse \Sequent \quad$}{}
  \end{enumerate}}
{berbentuk $!A \Sequent !A$} untuk sebarang !!{sentence} $!A$ dalam bahasa tersebut.
\end{defn}

!!^{derivation}s dalam kalkulus sekuen merupakan pohon-pohon sekuen
tertentu, dengan sekuen paling atas berupa sekuen awal; jika sebuah
sekuen berada di bawah satu atau dua sekuen lain, sekuen tersebut harus
mengikuti secara benar melalui suatu aturan inferensi. Aturan-aturan untuk
$\Log{LK}$ terbagi menjadi dua jenis utama: aturan \emph{logis} dan aturan
\emph{struktural}. Aturan logis dinamai menurut !!{main operator} dari
!!{sentence} yang memuat $!A$ dan/atau $!B$ dalam sekuen bawah. Setiap
aturan mempunyai dua versi: satu untuk menyimpulkan sekuen dengan
!!{sentence} yang memuat !!{operator} di sebelah kiri, dan satu lagi
dengan !!{sentence} tersebut di sebelah kanan.

\end{document}
